\section{Robustness and deployment stress posture}
\label{sec:robustness}

Robustness evidence is routed as stress-test evidence. The project includes
perturbation artifacts for sensor failure, observation noise, missingness,
heterogeneous costs, temporal shift, and candidate-count sensitivity. These
tests support bounded robustness discussion in the tested settings. They do not
support the phrase globally robust, and they do not replace held-out Stage15
evidence for main performance claims.

\begin{table*}[t]
\centering
\caption{Robustness evidence routing from Stage14 PeMS7\_228 stress tests.}
\label{tab:robustness-routing}
\scriptsize
\begin{tabularx}{\textwidth}{>{\raggedright\arraybackslash}p{0.22\textwidth}>{\raggedright\arraybackslash}p{0.22\textwidth}>{\centering\arraybackslash}p{0.12\textwidth}>{\centering\arraybackslash}p{0.12\textwidth}>{\centering\arraybackslash}p{0.10\textwidth}>{\raggedright\arraybackslash}X}
\toprule
Condition & Best layout & Best MAE & Layout rows & Splits & Claim route \\
\midrule
baseline & graph\_sampling\_laplacian & 1.2194 & 8 & 2 & stress-test only \\
block missing 12 steps & greedy\_d\_logdet & 1.1114 & 8 & 2 & stress-test only \\
cost proxy budget 45 & graph\_sampling\_laplacian & 1.2194 & 8 & 2 & stress-test only \\
observation noise 5\% & graph\_sampling\_laplacian & 1.2459 & 8 & 2 & stress-test only \\
random missing 10\% & graph\_sampling\_laplacian & 1.2143 & 8 & 2 & stress-test only \\
sensor failure 5\% & graph\_sampling\_laplacian & 1.2234 & 8 & 2 & stress-test only \\
sensor failure 10\% & graph\_sampling\_laplacian & 1.2308 & 8 & 2 & stress-test only \\
sensor failure 20\% & graph\_sampling\_laplacian & 1.2230 & 8 & 2 & stress-test only \\
chronological split & graph\_sampling\_laplacian & 1.1921 & 8 & 2 & stress-test only \\
\bottomrule
\end{tabularx}
\begin{flushleft}\footnotesize These rows are diagnostic Stage14 stress tests on PeMS7\_228. They route perturbation, cost, and temporal-shift evidence to bounded robustness discussion only; they do not replace the Stage15 TRACE-BiOpt dominance evidence.\end{flushleft}
\end{table*}


Table~\ref{tab:robustness-routing} makes this routing explicit. The rows are
not TRACE-BiOpt dominance rows; they are PeMS7\_228 Stage14 stress tests that
keep failure, noise, missingness, cost, and chronological-shift evidence visible
while preventing it from being promoted into a universal robustness claim. The
main TRACE-BiOpt result remains Table~\ref{tab:trace-biopt-dominance}.

\begin{table*}[t]
\centering
\caption{Stage14 PeMS7\_228 stress frontier over perturbation, cost, and temporal-shift slices. Lower MAE is better.}
\label{tab:trace-biopt-robustness-frontier}
\tiny
\setlength{\tabcolsep}{3pt}
\begin{tabular}{>{\raggedright\arraybackslash}p{0.18\textwidth}>{\raggedright\arraybackslash}p{0.16\textwidth}>{\raggedright\arraybackslash}p{0.12\textwidth}>{\centering\arraybackslash}p{0.09\textwidth}>{\raggedright\arraybackslash}p{0.16\textwidth}>{\centering\arraybackslash}p{0.08\textwidth}}
\toprule
Condition & Winner & Family & Best & Runner-up & Gap \\
\midrule
baseline & graph\_sampling & graph-spectral & 1.2194 & qr\_pod & 0.0297 \\
block missing 12 steps & greedy\_d\_logdet & logdet & 1.1114 & observability\_proxy & 0.0301 \\
cost proxy budget 45 & graph\_sampling & graph-spectral & 1.2194 & qr\_pod & 0.0297 \\
observation noise 5\% & graph\_sampling & graph-spectral & 1.2459 & qr\_pod & 0.0484 \\
random missing 10\% & graph\_sampling & graph-spectral & 1.2143 & multistart\_swap & 0.0406 \\
sensor failure 5\% & graph\_sampling & graph-spectral & 1.2234 & qr\_pod & 0.0081 \\
sensor failure 10\% & graph\_sampling & graph-spectral & 1.2308 & qr\_pod & 0.0284 \\
sensor failure 20\% & graph\_sampling & graph-spectral & 1.2230 & multistart\_swap & 0.0006 \\
chronological split & graph\_sampling & graph-spectral & 1.1921 & multistart\_swap & 0.0058 \\
\bottomrule
\end{tabular}
\begin{flushleft}\footnotesize These rows remain bounded Stage14 stress-test evidence only. They do not assert TRACE-BiOpt robustness under the same perturbations because the Stage14 artifact evaluates pre-TRACE-BiOpt reviewer-facing baselines on two split seeds; the table is used to summarize how the stress frontier moves across perturbation types.\end{flushleft}
\end{table*}


Table~\ref{tab:trace-biopt-robustness-frontier} adds the part that the routing
table alone hides: which baseline family actually sits on the Stage14 stress
frontier, and by how much. The graph-spectral surrogate
`graph\_sampling\_laplacian` wins 8 of the 9 perturbation slices, including all
sensor-failure, observation-noise, random-missing, and cost-proxy cases. The
only clear frontier shift is block missing 12 steps, where `greedy\_d\_logdet`
takes over. The narrowest gaps are chronological split (0.0058 MAE) and sensor
failure 20\% (0.0006 MAE), while observation noise 5\% and random missing 10\%
separate more clearly at 0.0484 and 0.0406 MAE. This is the right robustness
reading for the paper: the stress frontier does move across perturbation types,
but not in a way that supports a blanket robustness theorem for TRACE-BiOpt
itself. It supports bounded deployment discussion and motivates future
recoverability-driven robust extensions.

% Deployment stress posture table moved to supplementary material (SUBMISSION_EVIDENCE_INDEX.md)

The bounded robustness evidence pushes one step closer to deployment interpretation. The nominal and cost-proxy slices leave the same graph-spectral winner and gap, so the bounded stress check does not suggest a cost-driven posture change by itself. Diffuse corruption is also comparatively stable: observation noise and random missingness keep graph-spectral as the frontier winner with visibly wider gaps. The more cautionary slices are chronological drift and heavy failure. At chronological split, the frontier is already narrow at 0.0058 MAE, so temporal nonstationarity should trigger re-validation rather than a generic robustness claim. At 20\% sensor failure, the winner is still graph-spectral, but the gap shrinks to 0.0006 MAE, which is the clearest bounded stress sign that redundancy and fallback plans matter. The only tested family shift is block missing 12 steps, where logdet takes over, indicating that long contiguous outages can favor algebraic-observability structure over the nominal frontier winner. This is a deployment posture statement, not a TRACE-BiOpt robustness claim.
